\chapter{C++ for Latency III: The Compiler}\label{ch:ll:cpp-for-latency-iii-the-compiler}

In July 2009 a Linux kernel driver was found to contain the lines \texttt{struct sock *sk = tun->sk;} and, a little
further down, \texttt{if (!tun) return POLLERR;}. The check was meant to protect against a null pointer. The compiler had
deleted it: the pointer had already been dereferenced, a null dereference is \omterm{def:ll:cpp-for-latency-iii-the-compiler:ub}{undefined behaviour}, so the pointer could not
be null and the test was dead code. On a machine where address zero could be mapped, the deletion became an exploit. The
compiler did nothing wrong; it read the program's contract literally. This chapter is about that contract and about the
other levers the compiler offers a latency engineer: optimisation levels and target flags, profile-guided and \omterm{def:ll:cpp-for-latency-iii-the-compiler:lto}{link-time
optimisation}, and the habit that ties them together, reading the assembly and measuring every flag against a checksum.

\section{Undefined behaviour as an optimisation contract}

\begin{definition}[Undefined behaviour]\label{def:ll:cpp-for-latency-iii-the-compiler:ub}
\emph{Undefined behaviour}\index{undefined behaviour} is behaviour for which the C++ standard imposes no requirements:
signed integer overflow, a null or dangling dereference, an out-of-bounds access, a data race, reading an object through a
pointer of an unrelated type. The compiler may assume it never happens, and optimises every program as if it did not.
\end{definition}

The assumption is where much of the optimiser's power comes from. Because signed overflow cannot happen, \texttt{x + 1 > x}
is true, a loop counter cannot wrap, and a loop's trip count can be computed in advance; because dereferencing null cannot
happen, a pointer dereferenced once is non-null afterwards. The assembly in
\cref{lst:ll:cpp-for-latency-iii-the-compiler:plus,lst:ll:cpp-for-latency-iii-the-compiler:null} shows both on this book's
compiler: the comparison compiles to ``return 1'' even without optimisation, and the null check of the 2009 pattern is gone
at \texttt{-O2}.

\omcode{code/low-latency/08-cpp-for-latency-iii-the-compiler/cpp/asm/plus_one_greater_O2.s}{1}{2}{\texttt{plus\_one\_greater(int x)}, \texttt{return x + 1 > x;}, at -O2 (g++ 11.4): a constant.}\label{lst:ll:cpp-for-latency-iii-the-compiler:plus}
\omcode{code/low-latency/08-cpp-for-latency-iii-the-compiler/cpp/asm/first_then_check_O2.s}{1}{2}{\texttt{first\_then\_check(const int* p)}: dereference, then test for null; the test is gone.}\label{lst:ll:cpp-for-latency-iii-the-compiler:null}

For a trading system the lesson is double. \omterm{def:ll:cpp-for-latency-iii-the-compiler:ub}{Undefined behaviour} is not ``the machine does something odd''; it is a license
for the compiler to remove checks, reorder work and assume invariants that the code violates, and the result can change
with the compiler version or a flag. And the same contract is what makes C++ fast: the idioms that give the optimiser
facts (signed loop counters, no aliasing between unrelated types) are the ones that let it keep values in registers and
vectorise. The tools that catch violations run in test builds: the undefined-behaviour sanitiser
(\texttt{-fsanitize=undefined}) reports each overflow as it happens, the address sanitiser each invalid access.

\section{Aliasing}

\begin{definition}[Strict aliasing]\label{def:ll:cpp-for-latency-iii-the-compiler:alias}
Under the \emph{strict aliasing}\index{strict aliasing} rule an object may be accessed only through an lvalue of its own
type (or a closely related one: a signed or unsigned variant, a character type), so the compiler may assume that pointers
to unrelated types never refer to the same memory. GCC applies the rule from \texttt{-O2}.
\end{definition}

\omcode{code/low-latency/08-cpp-for-latency-iii-the-compiler/cpp/asm/alias_O2.s}{1}{4}{\texttt{alias(int* i, float* f)}: store 1 through \texttt{i}, 0 through \texttt{f}, return \texttt{*i}; the compiler returns 1 without reloading.}\label{lst:ll:cpp-for-latency-iii-the-compiler:alias}

In \cref{lst:ll:cpp-for-latency-iii-the-compiler:alias} the compiler returns 1 without reading \texttt{*i} again, because a
store through a \texttt{float*} cannot change an \texttt{int}. Called with two pointers to the same four bytes, the function
returns 1 where the memory holds 0. Decoders are where this bites: reading a price out of a byte buffer by casting the
buffer's address to \texttt{std::int32\_t*} is an aliasing violation and, for unaligned fields, also undefined. The correct
forms, \texttt{std::memcpy} into a local or C++20's \texttt{std::bit\_cast}, compile to a single load; the compiler knows
what they do (chapter~16 builds its codecs this way).

\section{Flags and target selection}

The optimisation level is a bundle of passes: \texttt{-O2} is the conventional release level, \texttt{-O3} adds more
aggressive \omterm{def:ll:cpp-for-latency-ii-compile-time:inline}{inlining}, cloning and vectorisation, and \texttt{-Ofast} adds \texttt{-ffast-math}. \texttt{-march=native}
lets the compiler use every instruction of the build machine (AVX2 on this laptop), which makes the binary unfit for older
machines; production builds name the target explicitly. \texttt{-ffast-math} allows the compiler to reorder floating-point
operations as if they were associative and to assume there are no NaNs or infinities.

\begin{definition}[Function multiversioning]\label{def:ll:cpp-for-latency-iii-the-compiler:fmv}
\emph{Function multiversioning}\index{function multiversioning} compiles one function several times for different
instruction sets and chooses a version at run time by the processor's features: GCC's \texttt{target\_clones(\textquotedbl{}avx2\textquotedbl{},
\textquotedbl{}default\textquotedbl{})} attribute builds both and dispatches when the program loads.
\end{definition}

\Cref{fig:ll:cpp-for-latency-iii-the-compiler:flags} measures the feed decoder of One Quant Book~1 with a per-message
handler in a second translation unit, under seven flag sets, and a floating-point sum under three. On the decoder,
\texttt{-O3} was more than three times faster than \texttt{-O2}, \omterm{def:ll:cpp-for-latency-iii-the-compiler:pgo}{profile-guided optimisation} alone two and a half times,
and the combination of everything nearly five times; \texttt{-march=native} and \omterm{def:ll:cpp-for-latency-iii-the-compiler:lto}{link-time optimisation} alone changed almost
nothing. On the float sum, \texttt{-ffast-math} was seven times faster and gave a different answer, because it summed in
another order: the benchmark's checksum caught it, and \texttt{firm.flagbench} does not rank a flag set whose output
differs.

\begin{omfigure}
\begin{tikzpicture}
\begin{axis}[name=dec,width=8.4cm,height=6cm,xbar,bar width=7pt,xmin=0,xmax=5,
  symbolic y coords={O0,O2,O2 native,O2 LTO,O2 PGO,O3,O3 native LTO PGO},ytick=data,enlarge y limits=0.1,
  xlabel={speed-up over -O2},title={\small feed decoder},tick label style={font=\footnotesize},label style={font=\small},
  grid=major,grid style={gray!25},nodes near coords,nodes near coords style={font=\scriptsize},
  every node near coord/.append style={/pgf/number format/fixed,/pgf/number format/precision=2},table/col sep=comma]
\addplot[fill=omDef!55,draw=omDef] table[y=name,x=speedup]{figdata/low-latency/08-cpp-for-latency-iii-the-compiler/measured_flags.csv};
\end{axis}
\begin{axis}[at={(dec.east)},anchor=west,xshift=1.9cm,width=5.8cm,height=6cm,xbar,bar width=9pt,xmin=0,xmax=9,
  symbolic y coords={O2,O3 native,O3 native fast-math},ytick=data,yticklabels={O2,O3 native,{O3 native\\fast-math}},
  y tick label style={align=right},enlarge y limits=0.3,xlabel={speed-up over -O2},title={\small float sum},
  tick label style={font=\footnotesize},label style={font=\small},grid=major,grid style={gray!25},nodes near coords,
  nodes near coords style={font=\scriptsize},every node near coord/.append style={/pgf/number format/fixed,/pgf/number format/precision=1},
  table/col sep=comma]
\addplot[fill=omThm!45,draw=omThm] table[y=name,x=speedup]{figdata/low-latency/08-cpp-for-latency-iii-the-compiler/measured_fastmath.csv};
\end{axis}
\end{tikzpicture}

\omcaption{Left: the decode benchmark (Book~1's 2\,020-message sample decoded 300 times, a per-message handler in another
translation unit) under seven flag sets; every set gave the same checksum. Right: summing $2^{20}$ floats; with
\texttt{-ffast-math} the sum changed. Best of five runs, g++ 11.4. Measured on a laptop (Intel Core Ultra 7 155H) under
WSL2, no isolated cores, pinned to one CPU. Data: \texttt{bench\_flags.py}.}\label{fig:ll:cpp-for-latency-iii-the-compiler:flags}
\end{omfigure}

\section{Profile-guided and link-time optimisation}

\begin{definition}[Profile-guided optimisation]\label{def:ll:cpp-for-latency-iii-the-compiler:pgo}
In \emph{profile-guided optimisation}\index{profile-guided optimisation} (PGO) the program is built with instrumentation,
run on a representative workload that records how often each branch goes each way and which values and calls are common,
and rebuilt with the profile, so that the compiler inlines, unrolls and lays out code according to measured behaviour
rather than guesses.
\end{definition}

\begin{definition}[Link-time optimisation]\label{def:ll:cpp-for-latency-iii-the-compiler:lto}
In \emph{link-time optimisation}\index{link-time optimisation} (LTO) the compiler stores its intermediate representation
in the object files and optimises the whole program again when linking, as if all translation units were one: functions
can be inlined across files.
\end{definition}

\begin{omfigure}
\begin{tikzpicture}[font=\footnotesize,
  st/.style={draw=omDef,fill=omDef!8,rounded corners=2pt,minimum height=0.75cm,align=center,text width=2.1cm},
  io/.style={draw=gray,fill=gray!10,rounded corners=2pt,minimum height=0.75cm,align=center,text width=1.9cm}]
\node[st] (a) at (0,0) {instrumented\\ build};
\node[st] (b) at (3.1,0) {train on a\\ replayed day};
\node[io] (c) at (6.0,0) {profile\\ (counts)};
\node[st] (d) at (8.9,0) {rebuild using\\ the profile};
\node[io] (e) at (11.6,0) {release\\ binary};
\foreach \x/\y in {a/b,b/c,c/d,d/e} \draw[-Stealth,thick] (\x) -- (\y);
\node[st] (f) at (0,-1.6) {each file to\\ intermediate code};
\node[st] (g) at (3.1,-1.6) {link: optimise\\ all files as one};
\node[io] (h) at (6.0,-1.6) {binary};
\draw[-Stealth,thick] (f) -- (g); \draw[-Stealth,thick] (g) -- (h);
\node[anchor=east,font=\footnotesize\bfseries] at (-1.3,0) {PGO};
\node[anchor=east,font=\footnotesize\bfseries] at (-1.3,-1.6) {LTO};
\end{tikzpicture}

\omcaption{\omterm{def:ll:cpp-for-latency-iii-the-compiler:pgo}{Profile-guided optimisation} builds twice around a training run, whose workload decides what the compiler
optimises for; \omterm{def:ll:cpp-for-latency-iii-the-compiler:lto}{link-time optimisation} defers the optimisation of the whole program to the link.}\label{fig:ll:cpp-for-latency-iii-the-compiler:pgo}
\end{omfigure}

The training run is the critical input of PGO. A profile taken on a quiet morning optimises the quiet-morning paths; a
trading system must be trained on a replay of a busy, representative day (chapter~23), and retrained when the workload
changes, or it will be optimised for the wrong branches. LTO's benefit depends on the code's structure: here the
per-message handler was small and its cost was dominated by the decoder's own calls, so LTO alone gained nothing. Both
lengthen the build and make the binary depend on more than the source; both belong in the release pipeline, measured by
the flag matrix on every release, not applied once and forgotten.

\section{Reading assembly}

The compiler's output answers questions no benchmark can: whether a function was inlined, a loop vectorised, a check
removed, a value kept in a register. \texttt{g++ -O2 -S -masm=intel} writes it for one file; \texttt{objdump -d
-M intel} disassembles a binary; online compiler explorers show it line by line. The book's listings of assembly are
generated by scripts from tested sources, and their tests check properties (no call, a constant return) with whatever
compiler is installed, never the exact bytes, which change with the compiler.

\begin{method}[Choosing release flags]\label{met:ll:cpp-for-latency-iii-the-compiler:flags}
\begin{enumerate}
\item Give the benchmark a checksum of its output and a realistic workload (a replayed day, not one message in a loop).
\item Build it under a matrix of flag sets, including profile-guided sets trained on a different day than the one measured.
\item Reject every set whose checksum differs from the reference build's, whatever its speed.
\item Rank the rest by measured percentiles over repeated runs (chapter~25), and read the assembly of the hot functions to
understand the winner.
\item Pin the chosen flags, compiler version and target in the build system; rerun the matrix on every compiler upgrade.
\end{enumerate}
\end{method}

\section{Tutorial: a flag matrix and three contracts}\label{tut:ll:cpp-for-latency-iii-the-compiler:tut}

\begin{tutorial}
\textbf{Goal.} See the optimiser act on \omterm{def:ll:cpp-for-latency-iii-the-compiler:ub}{undefined behaviour}, catch it with a sanitiser, and choose flags by measurement.
\textbf{End state:} the three assembly listings, a sanitiser report, and
\cref{fig:ll:cpp-for-latency-iii-the-compiler:flags}.
\begin{tutsteps}
\item \textbf{Three contracts.} \texttt{python ll\_ub.py} compiles \texttt{ll\_ub.cpp} to assembly at \texttt{-O0} and
\texttt{-O2} and writes the listings above.
\item \textbf{Catch an overflow.} \texttt{ll\_ub.ubsan(2147483647)} builds a call of \texttt{increment(x)}, \texttt{return
x + 1}, with \texttt{-fsanitize=undefined} and runs it: ``runtime error: signed integer overflow: 2147483647 + 1 cannot
be represented in type int''.
\item \textbf{Build with a profile.} \texttt{firm.flagbench} builds a profile-guided set in three steps: instrumented
build, training run, rebuild with the profile.
\omcode{code/firm/flagbench/firm_flagbench.py}{38}{51}{Building a flag set, profile-guided or not.}\label{lst:ll:cpp-for-latency-iii-the-compiler:build}
\item \textbf{Run the matrix} with \texttt{python bench\_flags.py}: seven sets on the decoder, three on the float sum, each
checked against the \texttt{-O2} checksum.
\end{tutsteps}
\textbf{What to change next.} Find which of the passes that \texttt{-O3} adds to \texttt{-O2} makes the decoder faster by
adding them one at a time to \texttt{-O2}; train the profile-guided build on only the first hundred messages of the sample
and measure again.
\end{tutorial}

\section{Build: the flag matrix}\label{bld:ll:cpp-for-latency-iii-the-compiler:flagbench}

\begin{build}
\bfield{Purpose} The firm's release flags are chosen, and re-chosen at every compiler upgrade, by measurement with a
correctness check, never by habit.
\bfield{Interface} Python \texttt{firm\_flagbench}: \texttt{FlagSet(name, flags, pgo)}, \texttt{build(sources, fs,
out\_dir, train\_args, includes)}, \texttt{run\_matrix(sources, flag\_sets, args, reference, repeats, cpus)} returning
rows with name, flags, checksum, time, ok and speed-up. The benchmark prints its checksum on its first line and its time
per unit on its second.
\bfield{Rules} A set whose checksum differs from the reference's is marked not ok and never recommended; profile-guided
sets are trained with the benchmark's own arguments (or a separate training workload); every run is pinned and repeated.
\bfield{Acceptance tests} \texttt{code/firm/flagbench/tests/}: a small program under \texttt{-O2}, \texttt{-O0} and a
profile-guided build (all ok) and under \texttt{-ffast-math} (not ok, its float sum differs).
\bfield{Stretch} Separate training and measurement workloads; report percentiles instead of the best time; bisect the
passes of \texttt{-O3} automatically.
\end{build}

\begin{omsources}
\item J.~Corbet, ``Fun with NULL pointers, part 1'', LWN.net, 20 July 2009.
\item C.~Lattner, ``What every C programmer should know about undefined behavior'', LLVM blog, 2011.
\item GCC 11.4 manual, ``Options that control optimization'' and ``Common function attributes''.
\end{omsources}

\section{Exercises}

\begin{exercise}[$\star$]\label{exo:ll:cpp-for-latency-iii-the-compiler:1}
From \cref{fig:ll:cpp-for-latency-iii-the-compiler:flags}, what does the decoder cost per message at \texttt{-O0},
\texttt{-O2} and with every optimisation, and what would each cost for 5 million messages a second?
\end{exercise}

\begin{exercise}[$\star$]\label{exo:ll:cpp-for-latency-iii-the-compiler:2}
Which of these are undefined: \texttt{INT\_MAX + 1}; \texttt{UINT\_MAX + 1u}; reading a \texttt{uint32\_t} from a byte
buffer with \texttt{std::memcpy}; reading it by casting the buffer's address to \texttt{uint32\_t*}?
\end{exercise}

\begin{exercise}[$\star$]\label{exo:ll:cpp-for-latency-iii-the-compiler:3}
Why does \texttt{plus\_one\_greater} compile to a constant even at \texttt{-O0}, and why can the sanitiser not catch it?
\end{exercise}

\begin{exercise}[$\star\star$]\label{exo:ll:cpp-for-latency-iii-the-compiler:4}
The exact sum of $1/k$ for $k = 1$ to $2^{20}$ is about 14.44. Which of the float sums of
\cref{fig:ll:cpp-for-latency-iii-the-compiler:flags} is closer to it, and does that make \texttt{-ffast-math} safe?
\end{exercise}

\begin{exercise}[$\star\star$]\label{exo:ll:cpp-for-latency-iii-the-compiler:5}
A decoder reads a 32-bit price with \texttt{*reinterpret\_cast<const std::uint32\_t*>(buf + 16)}. Name two rules it breaks
and write the correct line.
\end{exercise}

\begin{exercise}[$\star\star$]\label{exo:ll:cpp-for-latency-iii-the-compiler:6}
A profile-guided build was trained on a replay of a quiet day and deployed on the day of an index rebalance. What can go
wrong?
\end{exercise}

\begin{exercise}[$\star\star\star$]\label{exo:ll:cpp-for-latency-iii-the-compiler:7}
\emph{Coding.} Add a flag set \texttt{-O2 -fno-strict-aliasing} to the matrix and compile \texttt{alias} with it. What
changes in the assembly, and in the decoder's speed?
\end{exercise}

\begin{exercise}[$\star\star\star$]\label{exo:ll:cpp-for-latency-iii-the-compiler:8}
\emph{Find the flaw.} ``We compile with \texttt{-Ofast -march=native} on the build server; it was 30\% faster in our
benchmark.''
\end{exercise}

\section{Problem: The Fast Build That Was Wrong}

\begin{problem}[{Weekend problem --- a release flag review}]\label{pb:ll:cpp-for-latency-iii-the-compiler:1}
A team's release uses \texttt{-O2}. An engineer proposes \texttt{-O3 -march=native -flto} with \omterm{def:ll:cpp-for-latency-iii-the-compiler:pgo}{profile-guided
optimisation}, another \texttt{-Ofast}. The team runs \texttt{firm.flagbench}
(\cref{fig:ll:cpp-for-latency-iii-the-compiler:flags}).

\medskip\noindent\textbf{Part I --- The decoder.}
\begin{enumerate}
\item What speed-up over \texttt{-O2} does each flag set give?
\item Which single change gives most of the gain?
\item Why did \omterm{def:ll:cpp-for-latency-iii-the-compiler:lto}{link-time optimisation} alone give nothing here?
\item Did any set change the decoder's checksum?
\end{enumerate}

\medskip\noindent\textbf{Part II --- The float sum.}
\begin{enumerate}[resume]
\item What do the three sets give, in time and in value?
\item Why did \texttt{-ffast-math} change the value?
\item Which value is closer to the exact sum?
\item Why must the team still refuse \texttt{-ffast-math} for its pricing and risk code?
\end{enumerate}

\medskip\noindent\textbf{Part III --- \omterm{def:ll:cpp-for-latency-iii-the-compiler:ub}{Undefined behaviour}.}
\begin{enumerate}[resume]
\item The decoder reads fields by casting. Why might it work at \texttt{-O2} and break at \texttt{-O3}?
\item What does the undefined-behaviour sanitiser report for an overflowing increment?
\item In which builds should the sanitisers run?
\item What does the 2009 kernel bug teach a team that checks pointers after using them?
\end{enumerate}

\medskip\noindent\textbf{Part IV --- The decision.}
\begin{enumerate}[resume]
\item State the \emph{named result}: the speed-up of the combined optimisations over \texttt{-O2} on the decoder, and what
\texttt{-ffast-math} did to the float sum.
\item Which flags should the release use, and what must accompany them?
\item How should the profile's training day be chosen?
\item What does \texttt{-march=native} commit the firm to?
\item How would \omterm{def:ll:cpp-for-latency-iii-the-compiler:fmv}{function multiversioning} change that?
\item How often should the matrix be rerun?
\item Why check assembly properties rather than bytes in the tests?
\item In one sentence: what is a compiler flag worth without a checksum?
\end{enumerate}
\end{problem}

\section{Interview questions}

\begin{interviewq}[$\star$ \iqroles{developer}]\label{iq:ll:cpp-for-latency-iii-the-compiler:1}
What is \omterm{def:ll:cpp-for-latency-iii-the-compiler:ub}{undefined behaviour}, and why does it exist in C++?
\end{interviewq}

\begin{interviewq}[$\star\star$ \iqroles{developer}]\label{iq:ll:cpp-for-latency-iii-the-compiler:2}
What is the \omterm{def:ll:cpp-for-latency-iii-the-compiler:alias}{strict aliasing} rule? How do you read an integer from a network buffer correctly?
\end{interviewq}

\begin{interviewq}[$\star\star$ \iqroles{developer}]\label{iq:ll:cpp-for-latency-iii-the-compiler:3}
What does \omterm{def:ll:cpp-for-latency-iii-the-compiler:pgo}{profile-guided optimisation} do, and what are its risks?
\end{interviewq}

\begin{interviewq}[$\star\star$ \iqroles{developer}]\label{iq:ll:cpp-for-latency-iii-the-compiler:4}
Would you use \texttt{-ffast-math} in a trading system? Where, and where not?
\end{interviewq}

\begin{interviewq}[$\star\star$ \iqroles{developer}]\label{iq:ll:cpp-for-latency-iii-the-compiler:5}
A check \texttt{if (ptr == nullptr)} disappears from the optimised binary. Why might that be?
\end{interviewq}

\begin{interviewq}[$\star\star\star$ \iqroles{developer}]\label{iq:ll:cpp-for-latency-iii-the-compiler:6}
How would you decide, and keep deciding, the compiler flags of a latency-critical binary?
\end{interviewq}
